\subsection{The pure-trace core}
\label{pt:sec:environment}\label{lf:app:pure-trace}

Whenever Axiom~\textup{(A4)} is imposed below, fix one of its witnessing
outcomes and denote it by $o_\ast$.  Every channel between nonempty finite sets
$P:A\to\D(X)$ has the constant-payoff lift
\begin{equation}
K_P(o,x\mid a):=\1_{\{o=o_\ast\}}P(x\mid a).
\label{pt:eq:constant-payoff-lift}
\end{equation}
Throughout this constant-payoff phase, comparisons of pure record channels mean comparisons
of these lifts in the model already defined in the main text:
\[
q\succeq_Pq'\quad:\Longleftrightarrow\quad q\succeq_{K_P}q'.
\]
Likewise, $(q,P)\succeq(r,Q)$ denotes the main-text two-block comparison of
the lifts, and $P\sqcup Q$, processing, continuation lists, and
$P_1\triangleright(Q_1,\ldots,Q_k)$ denote the lifted main-text constructions.
The lift commutes with them up to tagged or singleton-record bijections,
neutral by Axiom~\textup{(A6)} in both directions.  These are only
abbreviations, not a new primitive relation or behavioural condition.

Two pieces of notation will be used repeatedly.  If
$S:A\to\Delta(B)$ is an action processor, let $S^qP:B\to\Delta(X)$ denote any
Bayesian pushforward satisfying
\[
 (qS)(b)(S^qP)(x\mid b)
 =\sum_a q(a)S(b\mid a)P(x\mid a).
\]
Its rows at zero-mass reports are arbitrary, exactly as in
\eqref{eq:action-processing}.  For channels $P:A\to\Delta(X)$ and
$Q:B\to\Delta(Y)$, their independent product is
\[
 (P\otimes Q)((x,y)\mid(a,b)):=P(x\mid a)Q(y\mid b).
\]

For $q\in\D(A)$ and $P:A\to\D(X)$, write
\[
m_{qP}(x):=\sum_a q(a)P(x\mid a),
\qquad
r_x^{qP}(a):=\frac{q(a)P(x\mid a)}{m_{qP}(x)}
\quad\text{when }m_{qP}(x)>0,
\]
and define the finite posterior law
\[
\mu_{qP}:=
\sum_{x:m_{qP}(x)>0}m_{qP}(x)\,\delta_{r_x^{qP}}.
\]
Let $\Id_A$ be full revelation and $U_A$ the one-record uninformative
channel.  With logarithms in the fixed base used in the main text, put
\[
I_{qP}(A;X)
:=\Sh(m_{qP})-\sum_a q(a)\Sh(P(\cdot\mid a))
=\Sh(q)-\int_{\D(A)}\Sh(r)\,d\mu_{qP}(r).
\]

\begin{proposition}[Pure-trace lemma]\label{pt:prop:main}
Suppose the preference family satisfies main-text Axioms
\textup{(A1)}, \textup{(A2)}, and \textup{(A4)--(A8)}.  Then, for every record channel between
nonempty finite sets
$P:A\to\D(X)$ and every $q,q'\in\D(A)$,
\[
q\succeq_Pq'
\quad\Longleftrightarrow\quad
I_{qP}(A;X)\ge I_{q'P}(A;X).
\]
The same numerical scale represents every block-supported comparison: for
every nonempty finite family $\{P_i:A_i\to\D(X_i)\}_{i\in J}$ of channels
between nonempty finite sets, distinct $i,j\in J$,
and priors $q_i\in\D(A_i)$ and $q_j\in\D(A_j)$,
\[
q_i^i\succeq_{\bigsqcup_{k\in J}P_k}q_j^j
\quad\Longleftrightarrow\quad
I_{q_iP_i}(A_i;X_i)\ge I_{q_jP_j}(A_j;X_j).
\]
\end{proposition}

Assume Proposition~\ref{pt:prop:main}'s hypotheses.  Every working property
below is the corresponding Axiom~\textup{(A1)}, \textup{(A2)}, or
\textup{(A4)--(A8)} applied to constant-payoff lifts.

\paragraph{Proof strategy.}
First, marked terminal laws give affine fixed-prior values, made exact under
relabelling and support transport.  Second, product and reached-branch
calibrations meet; their two-groupings argument eliminates the possible
product interaction and aligns every scale.  Third, the exact chain rule
produces a symmetric, continuous, strongly additive uncertainty functional,
identified as Shannon entropy by Faddeev's theorem.  The product-coefficient
and branch tangent-space lemmas may be read as technical interfaces.

\subsubsection{Fixed-prior affine values}

For a finite action set $A$, let $\mathsf X_A$ be the set of finitely
supported probability measures on $\Delta(A)$, and let
\[
 b(\mu):=\int_{\Delta(A)}s\,d\mu(s),
 \qquad
 \mathcal M_q:=\{\mu\in\mathsf X_A:b(\mu)=q\}.
\]
Thus $\mathcal M_q$ is the mixture set of finite Bayes-plausible posterior
laws with prior $q$, and $\mu_{qP}\in\mathcal M_q$ for every finite channel
$P$.  On the constant-payoff face, the marked terminal law from
\eqref{lf:eq:marked-law} is simply
\begin{equation}
 \eta_{K_P}=\delta_{o_\ast}\otimes\mu_{qP}.
 \label{pt:eq:pure-marked-law}
\end{equation}

Let $\delta_q\in\mathcal M_q$ be the posterior law of silence and let
\[
 \chi_q:=\sum_{a\in A}q(a)\delta_{e_a}\in\mathcal M_q
\]
be the posterior law of full revelation.

\begin{lemma}[Pure affine fibres]
\label{pt:lem:pure-affine}
\label{pt:lem:blockcoh}
\label{pt:lem:blackwell}\label{pt:lem:plsuff}
\label{pt:lem:publiccoin}
\label{pt:lem:postsep}
\label{pt:lem:convnorm}
\emph{Global order facts.}
Pure prior--channel pairs form a weak order with replacement, unaffected by
adding or deleting other blocks.  For every channel $P$ and priors $q,r$ on
its action set,
\begin{equation}
 q\succeq_P r
 \quad\Longleftrightarrow\quad
 (q,P)\succeq(r,P),
 \label{pt:eq:within-to-pair}
\end{equation}
and on fixed alphabets the pair-order graph is closed under entrywise
convergence of channels and priors.  \emph{Fixed full-support fibre.}
Fix a full-support $q\in\Delta(A)$.
Then the following statements hold.

\begin{enumerate}[label=\textup{(\alph*)},leftmargin=*]
\item If $\mu_{qP}=\mu_{qQ}$, record processors exist in both directions,
  so $(q,P)\sim(q,Q)$ and either channel may replace the other in a block
  comparison.  Conversely, every
  $\mu=\sum_{k=1}^n\lambda_k\delta_{s_k}\in\mathcal M_q$, written with
  $\lambda_k>0$, is implemented by
  \begin{equation}
   C_\mu(k\mid a):=\frac{\lambda_ks_k(a)}{q(a)},
   \qquad \mu_{qC_\mu}=\mu.
   \label{pt:eq:canonical-implementation}
  \end{equation}
  Hence the fixed-$q$ order descends to all of $\mathcal M_q$.

\item For publicly tagged mixtures,
  \begin{equation}
   \mu_{q,tP\oplus(1-t)R}
   =t\mu_{qP}+(1-t)\mu_{qR},
   \label{pt:eq:public-posterior-mixture}
  \end{equation}
  and, for $0<t<1$,
  \[
   (q,P)\succeq(q,Q)
   \quad\Longleftrightarrow\quad
   (q,tP\oplus(1-t)R)\succeq(q,tQ\oplus(1-t)R).
  \]

\item There is an affine $F_q:\mathcal M_q\to\mathbb R$ representing the
  fixed-prior order:
  \begin{equation}
   (q,P)\succeq(q,Q)
   \quad\Longleftrightarrow\quad
   F_q(\mu_{qP})\ge F_q(\mu_{qQ}).
   \label{pt:eq:fibre-representation}
  \end{equation}
  It may be normalized by
\begin{equation}
 F_q(\delta_q)=0,
 \qquad
 F_q(\mu)\ge0\quad(\mu\in\mathcal M_q),
 \label{pt:eq:zero-minimum}
\end{equation}
and, whenever $|A|\ge2$,
\begin{equation}
 F_q(\chi_q)=1.
 \label{pt:eq:full-revelation-anchor}
\end{equation}
These two anchors select a unique affine representative on every
nondegenerate full-support fibre.  On a singleton fibre we retain
$F_q\equiv0$.
\end{enumerate}
\end{lemma}

\begin{proof}
The order, replacement, and closed-graph assertions are respectively
Lemma~\ref{lf:lem:bookkeeping}\textup{(a)}, the duplication clause of
Axiom~\textup{(A5)}, and Axiom~\textup{(A2)}, all specialized to the lifts.

For part \textup{(a)}, equation~\eqref{pt:eq:pure-marked-law} turns equality
of posterior laws into equality of marked terminal laws.  The matching
processors in Lemma~\ref{lf:lem:terminal-affine}\textup{(a)} therefore give
record processors in both directions, including arbitrary completions on
zero-mass rows.  Replacement gives the quotient order.  Conversely, Bayes
plausibility gives $\sum_k\lambda_ks_k=q$, so the rows in
\eqref{pt:eq:canonical-implementation} sum to one and Bayes' formula gives
$\mu_{qC_\mu}=\mu$.

For part \textup{(b)}, Bayes' formula gives
\eqref{pt:eq:public-posterior-mixture}.  On constant-payoff lifts this is the
public-mixture identity~\eqref{lf:eq:public-mixture}, and the independence
argument in Lemma~\ref{lf:lem:terminal-affine}\textup{(b)} gives the
biconditional.  Tagged disjoint unions handle unequal record alphabets, with
part \textup{(a)} making the padding neutral.

For part \textup{(c)}, restrict the affine representative from
Lemma~\ref{lf:lem:terminal-affine}\textup{(b)} to laws
$\delta_{o_\ast}\otimes\mu$ and use parts \textup{(a)--(b)} and canonical
implementability.  If $|A|\ge2$, Lemma~\ref{lem:relevance-bridge}\textup{(b)}
makes the order nonconstant, so positive-affine uniqueness applies.  Deleting
the record garbles every channel to silence; hence $\delta_q$ is a minimum.
Translate its value to zero and divide by the positive value of $\chi_q$.
This gives the two anchors and nonnegativity.  If $|A|=1$, Bayes plausibility
makes $\mathcal M_q=\{\delta_q\}$.
\end{proof}

The only continuity used here is the finite-alphabet segment continuity
already verified in Lemma~\ref{lf:lem:terminal-affine}: implement the finitely
many laws on one tagged record alphabet and apply Axiom~\textup{(A2)} along
their mixtures.  No topology on laws with growing supports, and no pointwise
integrand on $\Delta(A)$, is being assumed.

\subsubsection{Exact transport across priors and support faces}

Action processing handles relabelling, undetectable distinctions, and
restriction to the positive support in one step.

For a bijection $\pi:A\to A'$, write $q\pi$ for the transported prior and
$\pi_\#\mu$ for the posterior law obtained by transporting every atom of
$\mu$.  If $B=\operatorname{supp}(q)$, write $q_B$ for the full-support
restriction of $q$ to $B$, $P_B$ for the restricted channel, and
$\pi_B\mu$ for the law obtained by restricting every posterior atom to $B$.

\begin{lemma}[Relabelling, undetectable distinctions, and support transport]
\label{pt:lem:transport}\label{pt:lem:actionbase}\label{pt:lem:neutrality}\label{pt:lem:supprestrict}
The canonically normalized fibre values have the following properties.

\begin{enumerate}[label=\textup{(\alph*)},leftmargin=*]
\item \emph{Exact relabelling.}
If $q$ has full support and $\pi:A\to A'$ is a bijection, then
\begin{equation}
 F_{q\pi}^{A'}(\pi_\#\mu)=F_q^A(\mu)
 \qquad(\mu\in\mathcal M_q).
 \label{pt:eq:exact-relabel}
\end{equation}

\item \emph{Undetectable distinctions.}
Let $S:A\to B$ be an onto deterministic map, let $q$ have full support, and
suppose that $P:A\to\Delta(X)$ is constant on every fibre of $S$.  Then
\begin{equation}
 (q,P)\sim(qS,S^qP).
 \label{pt:eq:undetectable}
\end{equation}
In particular, if $G_S$ reports only $S(a)$, then
$(q,G_S)\sim(qS,\Id_B)$.

\item \emph{Support restriction.}
For an arbitrary prior $q\in\Delta(A)$ with nonempty support
$B:=\operatorname{supp}(q)$ and every pair of finite channels
$P:A\to\Delta(X)$ and $Q:A\to\Delta(Y)$,
\begin{equation}
 (q,P)\sim(q_B,P_B).
 \label{pt:eq:support-indifference}
\end{equation}
Consequently comparisons at $q$ depend only on the rows indexed by $B$, and
\begin{equation}
 (q,P)\succeq(q,Q)
 \quad\Longleftrightarrow\quad
 (q_B,P_B)\succeq(q_B,Q_B).
 \label{pt:eq:support-comparison}
\end{equation}
Every atom of a law $\mu\in\mathcal M_q$ is supported on $B$, and the
definition
\begin{equation}
 F_q^A(\mu):=F_{q_B}^{B}(\pi_B\mu)
 \label{pt:eq:boundary-representative}
\end{equation}
represents all continuation comparisons at the boundary prior $q$.
\end{enumerate}
\end{lemma}

\begin{proof}
\emph{Part (a).}
Axiom~\textup{(A7)} applied to $\pi$ and $\pi^{-1}$ identifies the two
orders.  Hence $\mu\mapsto F_{q\pi}^{A'}(\pi_\#\mu)$ and $F_q^A$ are affine
representatives of the same order.  Relabelling preserves silence and full
revelation, so the two anchors and affine uniqueness give
\eqref{pt:eq:exact-relabel}; singleton values vanish.

\emph{Part (b).}
One direction of \eqref{pt:eq:undetectable} is Axiom~\textup{(A7)}.  For the
reverse, use the Bayesian refinement kernel
\[
 R(a\mid b):=
 \begin{cases}
 q(a)/(qS)(b),&S(a)=b,\\
 0,&S(a)\ne b.
 \end{cases}
\]
It satisfies $(qS)R=q$, and fibre constancy means that pushing $S^qP$ back
through $R$ recovers $P$.  A second application of Axiom~\textup{(A7)} gives
the reverse comparison.  Taking $P=G_S$ gives the final assertion.

\emph{Part (c).}
Lemma~\ref{lf:lem:bookkeeping}\textup{(b)} gives
\eqref{pt:eq:support-indifference}, and replacement gives
\eqref{pt:eq:support-comparison}.  If
$\mu=\sum_k\lambda_k\delta_{s_k}\in\mathcal M_q$ with $\lambda_k>0$, then
$0=q(a)=\sum_k\lambda_ks_k(a)$ off $B$; nonnegativity puts every $s_k$ on
$\Delta(B)$.  Thus $\pi_B\mu\in\mathcal M_{q_B}$, and the full-support
representation on $B$ proves~\eqref{pt:eq:boundary-representative}.
Compatibility with relabelling follows from part~\textup{(a)}.
\end{proof}

Henceforth action-set superscripts are suppressed, and boundary subscripts
mean transport from the positive support.

\paragraph{Normalization ledger.}
$F_q$ has the silence/full-revelation normalization; product calibration gives
$G_q=c_qF_q$ with a possible interaction $\kappa$; branch calibration gives
$V_q=G_q/a_q$.  Compatibility will make $a_q$ universal and $\kappa=0$.

\subsubsection{Scale coherence}

\paragraph{Product calibration and a common comparison scale.}

Write $P\succeq_qQ$ for $(q,P)\succeq(q,Q)$, and similarly for strict
preference.  For a channel $P:A\to\Delta(X)$, let $\widetilde P$ denote its
lift to $A\times B$, constant in the $B$-coordinate; for
$R:B\to\Delta(Y)$, let $\widetilde R$ denote the symmetric lift, constant in
the $A$-coordinate.
By Lemma~\ref{pt:lem:plsuff}, null records and record permutations are neutral.

\begin{lemma}[Cancellation of a common independent background]
\label{pt:lem:background}
Let $p\in\Delta(A)$ and $r\in\Delta(B)$ have full support.  For channels
$P,Q$ on $A$ and a channel $R$ on $B$,
\begin{equation}
 P\otimes R\succeq_{p\otimes r}Q\otimes R
 \quad\Longleftrightarrow\quad
 P\succeq_p Q.
\tag{BI}\label{pt:eq:backgroundinert}
\end{equation}
The symmetric assertion holds with the two coordinates interchanged.  The
equivalence, including its strict form, remains valid when a posterior reached
in the background coordinate lies on a proper support face.
\end{lemma}

\begin{proof}
Pad $P$ and $Q$ by null records so that they have the same record alphabet.
Regard the lift $\widetilde R$ of $R$ as a first-stage experiment on
$A\times B$, followed on every reached branch by $\widetilde P$, respectively
$\widetilde Q$.  Up to the harmless permutation of the two record
coordinates, the resulting compounds are $P\otimes R$ and $Q\otimes R$.

If $y$ is reached under $R$, its posterior on $A\times B$ is
\[
 p\otimes r_y,
 \qquad r_y:=r_y^{\,r,R}.
\]
The coordinate projection $A\times\operatorname{supp}(r_y)\to A$ and
Lemma~\ref{pt:lem:transport}\textup{(b)--(c)} give
\[
 (p\otimes r_y,\widetilde P)\sim(p,P),
 \qquad
 (p\otimes r_y,\widetilde Q)\sim(p,Q).
\tag{BI$'$}\label{pt:eq:bg-branch-transport}
\]
Thus every reached branch has the same weak comparison as $P$ versus $Q$ at
$p$.  If $P\succeq_pQ$, Lemma~\ref{pt:lem:finite-branch-extension} applied to
these branches proves the forward implication in~\eqref{pt:eq:backgroundinert}.

For the converse, suppose that $P\not\succeq_pQ$.  Completeness gives
$Q\succ_pP$.  By~\eqref{pt:eq:bg-branch-transport}, the reversed comparison is
strict on every reached branch; at least one such branch has positive mass.
The strict part of Lemma~\ref{pt:lem:finite-branch-extension} therefore gives
$Q\otimes R\succ_{p\otimes r}P\otimes R$, contradicting
$P\otimes R\succeq_{p\otimes r}Q\otimes R$.  This proves the converse and
also records why weak branch monotonicity alone would not suffice.  Interchange
the coordinates for the symmetric assertion.
\end{proof}

For finite posterior laws define
\[
 \left(\sum_i m_i\delta_{u_i}\right)\otimes
 \left(\sum_j n_j\delta_{v_j}\right)
 :=\sum_{i,j}m_i n_j\delta_{u_i\otimes v_j}.
\]
Bayes' rule gives
\begin{equation}
 \mu_{p\otimes r,P\otimes R}=\mu_{pP}\otimes\mu_{rR}.
\tag{PL}\label{pt:eq:productlaw}
\end{equation}

\begin{lemma}[Product gauge]
\label{pt:lem:coherentnorm}\label{pt:cor:permutationinvariance}
There are positive rescalings $G_q=c_qF_q$ of the zero-normalised affine
representatives and a single number $\kappa\in\mathbb R$ such that
\begin{equation}
 G_{p\otimes r}(\mu\otimes\nu)
 =G_p(\mu)+G_r(\nu)+\kappa G_p(\mu)G_r(\nu)
\tag{QA}\label{pt:eq:QA}
\end{equation}
for all full-support priors $p,r$ and all
$\mu\in\mathcal M_p$, $\nu\in\mathcal M_r$.  These rescalings can be chosen
exactly invariant under bijections.  If a factor is a singleton, its
representative is zero and the same formula holds.  Finally,
\begin{equation}
 1+\kappa G_r(\nu)>0
\tag{POS}\label{pt:eq:QAslope}
\end{equation}
whenever the $r$-fibre is nonconstant.
\end{lemma}

\begin{proof}
We give the algebra because both the common value of $\kappa$ and the strict
inequality~\eqref{pt:eq:QAslope} are used later.  Initially retain the notation
$F_q$.  For nondegenerate $p,r$, put
\[
 L_{p,r}(\mu,\nu):=F_{p\otimes r}(\mu\otimes\nu),
 \qquad x:=F_p(\mu),\quad y:=F_r(\nu).
\]
The map $L_{p,r}$ is separately affine, and
Lemma~\ref{pt:lem:background} says that each slice represents the appropriate
factor order.  We record the short uniqueness argument.  Fixing $\nu$ gives
$L_{p,r}=\alpha(\nu)x+\gamma(\nu)$ with $\alpha(\nu)>0$.  At
$\mu=\delta_p$, affine uniqueness gives
$\gamma(\nu)=B_{p,r}y$ for some $B_{p,r}>0$.  At any $\bar\mu$ with
$\bar x:=F_p(\bar\mu)\ne0$, the other slice has the form
$A_{p,r}\bar x+D_{p,r}y$ with $A_{p,r}>0$.  Comparing the two expressions
gives $\alpha(\nu)=A_{p,r}+C_{p,r}y$, where
$C_{p,r}:=(D_{p,r}-B_{p,r})/\bar x$.  Hence
\begin{equation}
 L_{p,r}(\mu,\nu)
 =A_{p,r}x+B_{p,r}y+C_{p,r}xy.
\label{pt:eq:pairbilinear}
\end{equation}
The symmetric calculation and strict slice cancellation give
\begin{equation}
 A_{p,r}+C_{p,r}y>0,
 \qquad B_{p,r}+C_{p,r}x>0
\label{pt:eq:pair-slopes}
\end{equation}
throughout the value ranges.

Exact relabelling makes the evaluations under the two associations of a
triple product equal.  Comparing the coefficients of $x,y,z$ gives
\begin{align}
 A_{p\otimes r,s}A_{p,r}&=A_{p,r\otimes s},\label{pt:eq:Acoc1}\\
 A_{p\otimes r,s}B_{p,r}&=B_{p,r\otimes s}A_{r,s},\label{pt:eq:Acoc2}\\
 B_{p\otimes r,s}&=B_{p,r\otimes s}B_{r,s}.
\label{pt:eq:Acoc3}
\end{align}
The coordinate swap similarly gives
\[
 B_{r,p}=A_{p,r},\qquad A_{r,p}=B_{p,r},\qquad C_{r,p}=C_{p,r}.
\tag{SW}\label{pt:eq:coefficient-swap}
\]
This coefficient comparison is legitimate: the laws
$(1-t)\delta_p+t\chi_p$ make $x$ range over a nondegenerate interval, and
the same holds for $y$ and $z$.
Let $\rho(p,r):=B_{p,r}/A_{p,r}$.  Dividing
\eqref{pt:eq:Acoc2} by~\eqref{pt:eq:Acoc1} yields
\[
 \rho(p,r)=\rho(p,r\otimes s)A_{r,s}.
\]
Fix a nondegenerate reference prior $q_0$ and set
$g(r):=\rho(q_0,r)$.  Then
\begin{equation}
 A_{r,s}=\frac{g(r)}{g(r\otimes s)}.
\label{pt:eq:gaugeA}
\end{equation}
The uniqueness of the coefficients in~\eqref{pt:eq:pairbilinear}, applied
after product relabellings, shows that $g$ is invariant under bijections.
Consequently the positive rescaling $G_p:=g(p)F_p$ preserves exact
relabelling.  Equation~\eqref{pt:eq:gaugeA} makes the new first coefficient
equal to one, and~\eqref{pt:eq:coefficient-swap} makes the second coefficient
equal to one as well.  We henceforth use $G$ for these gauged representatives.

Write the remaining coefficient as $\kappa_{p,r}$.  Associativity of
\[
 x\oplus_{p,r}y:=x+y+\kappa_{p,r}xy
\]
and comparison of the coefficients of $xy,xz,yz$ give
\[
 \kappa_{p,r}=\kappa_{p,r\otimes s},\qquad
 \kappa_{p\otimes r,s}=\kappa_{p,r\otimes s},\qquad
 \kappa_{p\otimes r,s}=\kappa_{r,s}.
\tag{KA}\label{pt:eq:kappa-assoc}
\]
(The $xyz$ coefficient gives the product of the corresponding two sides and
adds no restriction.)  Taking $s=q_0$ in~\eqref{pt:eq:kappa-assoc} shows that
$\kappa_{p,r}=\kappa_{r,q_0}$, so it is independent of $p$; the swap in
\eqref{pt:eq:coefficient-swap} makes it symmetric, hence independent of $r$.
Call the common value $\kappa$.  After the gauge,
\eqref{pt:eq:pair-slopes} is exactly
$1+\kappa G_r(\nu)>0$ and its symmetric counterpart.

For a singleton prior set $G\equiv0$; the canonical bijections
$A\cong A\times\{*\}$ give~\eqref{pt:eq:QA}.  For a boundary prior, define
$G_r$ by exact transport from its positive-probability support, as in
Lemma~\ref{pt:lem:supprestrict}.  This convention is relabelling invariant.
\end{proof}

\begin{lemma}[Cross-prior block representation]
\label{pt:lem:blockbridge}
For arbitrary finite priors $q\in\Delta(A)$ and $r\in\Delta(B)$ and channels
$P,Q$ on the respective action sets,
\begin{equation}
 q^0\succeq_{P\sqcup Q}r^1
 \quad\Longleftrightarrow\quad
 G_q(\mu_{qP})\ge G_r(\mu_{rQ}).
\tag{BC}\label{pt:eq:blockcomp}
\end{equation}
\end{lemma}

\begin{proof}
First suppose that $q,r$ have full support.  By~\eqref{pt:eq:QA} and
$G_q(\delta_q)=G_r(\delta_r)=0$,
\[
 G_{q\otimes r}(\mu_{qP}\otimes\delta_r)=G_q(\mu_{qP}),
 \qquad
 G_{q\otimes r}(\delta_q\otimes\mu_{rQ})=G_r(\mu_{rQ}).
\]
The two laws on the left lie in the same fibre and are implemented by
$P\otimes U_B$ and $U_A\otimes Q$, so the affine fibre representation
compares them by the displayed numbers.

Lemma~\ref{lf:lem:bookkeeping}\textup{(c)} removes the independent dummy
coordinates, giving
$(q\otimes r,P\otimes U_B)\sim(q,P)$ and
$(q\otimes r,U_A\otimes Q)\sim(r,Q)$.  Replacement proves
\eqref{pt:eq:blockcomp}.  Arbitrary priors follow by support restriction;
singleton supports use the zero convention.
\end{proof}

\paragraph{Branch calibration and coherent face scales.}

Changing a continuation on one reached branch changes the posterior law by a
signed measure with zero total mass and zero barycentre.  The following space
is exactly the space of such feasible directions.
For a nonempty finite set $A$, let
\[
 W_A:=\left\{\eta:
 \begin{array}{l}
 \eta\text{ is a finitely supported signed measure on }\Delta(A),\\[-2pt]
 \eta(\Delta(A))=0,\quad \displaystyle\int p\,d\eta(p)=0
 \end{array}\right\}.
\]
If $\iota:B\hookrightarrow A$ is an injection, let
$i_\iota:W_B\to W_A$ push every posterior forward along $\iota$; for a
literal subset $B\subseteq A$, write $i_B$ for the inclusion case.  Let
$L_q^A$ be the linear part of the affine functional $G_q$ on the affine hull
of $\mathcal M_q$.

\begin{lemma}[Branch aggregation]
\label{pt:lem:branchagg}
Let $q$ have full support on $A$.  For every $r\in\Delta(A)$ with
$B:=\operatorname{supp}(r)$ and $|B|\ge2$, there is a unique
$\beta_A(q,r)>0$ satisfying
\begin{equation}
 L_q^A\circ i_B=\beta_A(q,r)L_{r_B}^B
 \quad\hbox{on }W_B.
\label{pt:eq:branch-linear}
\end{equation}
It depends only on $q,r$ and the face inclusion.  If $k\ge1$ and
$P_1:A\to\Delta(\{1,\ldots,k\})$, define
\begin{equation}
 m_i:=\sum_aq(a)P_1(i\mid a),\qquad
 r_i(a):=\frac{q(a)P_1(i\mid a)}{m_i}\quad(m_i>0).
\label{pt:eq:branch-notation}
\end{equation}
Then every list of finite continuation channels $Q_1,\ldots,Q_k$ on $A$
satisfies
\begin{equation}
 G_q\bigl(\mu_{q,P_1\triangleright(Q_1,\ldots,Q_k)}\bigr)
 =G_q(\mu_{qP_1})+
 \sum_{i:m_i>0}m_i\beta_A(q,r_i)
 G_{r_i}(\mu_{r_i,Q_i}).
\tag{BA}\label{pt:eq:branchaggregation}
\end{equation}
For singleton $r_i$, the continuation value is zero and the corresponding
coefficient may be fixed by any positive convention.
\end{lemma}

\begin{proof}
We first record the common tangent space.  If $q$ has full support, then
\begin{equation}
 \operatorname{span}\{\mu-\nu:\mu,\nu\in\mathcal M_q\}=W_A.
\label{pt:eq:tangent-space}
\end{equation}
Only the reverse inclusion needs proof.  For $0\ne\eta\in W_A$, write its
Jordan decomposition as $\eta=\eta^+-\eta^-$, where both positive parts have
the same mass $c>0$ and, after division by $c$, the same barycentre $b$.
Choose $\varepsilon>0$ small enough that
$\tau=(q-\varepsilon b)/(1-\varepsilon)$ belongs to $\Delta(A)$.  Then
\[
 \mu^\pm:=\frac{\varepsilon}{c}\eta^\pm
          +(1-\varepsilon)\delta_\tau
 \in\mathcal M_q,
 \qquad
 \eta=\frac c\varepsilon(\mu^+-\mu^-),
\]
which proves~\eqref{pt:eq:tangent-space}.  The same argument intrinsically on
a face gives the space $W_B$.

Fix a nondegenerate $r$ supported on $B$.  It is reachable from $q$: choose
\[
 0<\varepsilon<\min\left\{1,
       \min_{b\in B}\frac{q(b)}{r(b)}\right\}
\]
and use a two-record experiment whose first record has conditional probability
$\varepsilon r(a)/q(a)$.  That record has mass $\varepsilon$ and posterior
$r$.

Implement $\nu,\nu'\in\mathcal M_{r_B}$ by channels on $B$, extend their
off-support rows arbitrarily to $A$, and pad their record alphabets if
necessary.  Vary only the continuation on the branch reaching $r$ from
$\nu$ to $\nu'$.  The aggregate posterior law changes by
$\varepsilon i_B(\nu-\nu')$.  Corollary~\ref{pt:cor:one-branch-a8} and
support transport give
\[
 \operatorname{sgn}L_q^Ai_B(\nu-\nu')
 =\operatorname{sgn}L_{r_B}^B(\nu-\nu').
\]
Indeed, strict preference handles the positive case, swapping continuations
handles the negative case, and indifference gives both inequalities in the
zero case.  By~\eqref{pt:eq:tangent-space}, this sign identity extends by
positive scaling to $W_B$.  Positive trace orientation makes
\[
 L_{r_B}^B(\chi_{r_B}-\delta_{r_B})
 =G_{r_B}(\chi_{r_B})>0,
\]
so $L_{r_B}^B\ne0$.  Two nonzero linear functionals with the same positive
half-space have the same kernel and are positive scalar multiples; this gives
the unique coefficient in~\eqref{pt:eq:branch-linear}.  Since both linear
functionals in that equation are fixed before a reaching experiment is
chosen, the coefficient is path-independent.

For the aggregation formula, put $B_i:=\operatorname{supp}(r_i)$ and
\[
 \nu_i:=\pi_{B_i}\mu_{r_i,Q_i}\in\mathcal M_{(r_i)_{B_i}}.
\]
The compound law differs from the first-stage law by
\[
 \mu_{q,P_1\triangleright(Q_1,\ldots,Q_k)}-\mu_{qP_1}
 =\sum_{i:m_i>0}m_i\,i_{B_i}
   (\nu_i-\delta_{(r_i)_{B_i}}).
\]
Apply $L_q^A$, use~\eqref{pt:eq:branch-linear}, and use
$G_{r_i}(\delta_{r_i})=0$.  Singleton summands vanish.  This is
\eqref{pt:eq:branchaggregation}.
\end{proof}

For full-support $q\in\Delta(A)$, an injection
$\iota:B\hookrightarrow A$, and full-support $r\in\Delta(B)$, define
\[
 \beta_\iota(q,r):=\beta_A(q,\iota_\#r).
\]
The image posterior $\iota_\#r$ is reachable by the two-record construction
in the proof, and its support inclusion is precisely $i_\iota$.  Thus
\[
 L_q^A i_\iota=\beta_\iota(q,r)L_r^B
 \quad\hbox{on }W_B.
\]

\begin{lemma}[Cocycle, face coherence, and the chain gauge]
\label{pt:lem:chain}
There is a relabelling-invariant family of positive scales
\[
\{a_q^A: A\text{ finite},\ |A|\ge2,\ q\in\Delta(A)\text{ full support}\}
\]
for the coefficients of Lemma~\ref{pt:lem:branchagg}, such that
\begin{equation}
 \beta_\iota(q,r)=\frac{a_q^A}{a_r^B}
\tag{FS}\label{pt:eq:facescale}
\end{equation}
for every injection $\iota:B\hookrightarrow A$ with $|A|,|B|\ge2$ and every
full-support $q\in\Delta(A)$ and $r\in\Delta(B)$; here
$\beta_\iota(q,r)$ is the coefficient for reaching the posterior
$\iota_\#r$ on the face $\iota(B)$.  For each prior in this family, put
\[
 V_q:=\frac{G_q}{a_q}
\].
These representatives satisfy the following exact chain rule.  For every
finite $A$ with $|A|\ge2$, every full-support $q\in\Delta(A)$, every $k\ge1$,
every $P_1:A\to\Delta(\{1,\ldots,k\})$, and every list of finite continuation
channels $Q_1,\ldots,Q_k$ on $A$, use the branch quantities in
\eqref{pt:eq:branch-notation}.  Then
\begin{equation}
 V_q\bigl(\mu_{q,P_1\triangleright(Q_1,\ldots,Q_k)}\bigr)
 =V_q(\mu_{qP_1})
  +\sum_{i:m_i>0}m_iV_{r_i}(\mu_{r_i,Q_i}).
\tag{CR}\label{pt:eq:chainrule}
\makeatletter
\g@addto@macro\df@tag{\ltx@label{pt:eq:final-chain}}
\makeatother
\end{equation}
Singleton continuation terms are zero.
\end{lemma}

\begin{proof}
There are two normalization tasks: first solve the branch cocycle within each
fixed action-set dimension, then remove the residual scalar attached to
embedding an $m$-point face in an $n$-point simplex.

We first prove the cocycle before choosing scales.  For nested injections
$C\xhookrightarrow{\jmath}B\xhookrightarrow{\iota}A$ and full-support
priors $q,r,s$ on $A,B,C$, respectively, the defining identities on $W_C$
are
\[
 L_q^A i_{\iota\circ\jmath}
 =\beta_\iota(q,r)L_r^B i_\jmath
 =\beta_\iota(q,r)\beta_\jmath(r,s)L_s^C.
\]
The direct coefficient from $q$ to $s$ is unique because $L_s^C\ne0$ by
Axiom~\textup{(A4)}.  Therefore
\begin{equation}
 \beta_{\iota\circ\jmath}(q,s)
 =\beta_\iota(q,r)\beta_\jmath(r,s).
\label{pt:eq:branch-cocycle}
\end{equation}
This single identity covers full-support paths, full-to-face paths, and
nested support faces.

For a fixed nondegenerate $A$, apply~\eqref{pt:eq:branch-cocycle} to identity
embeddings.  Fix a full-support basepoint $q_A$ and put
\[
 a_q^A:=\beta_{\mathrm{id}_A}(q,q_A).
\]
Then
\[
 \beta_{\mathrm{id}_A}(q,r)=\frac{a_q^A}{a_r^A}.
\tag{WS}\label{pt:eq:within-scale}
\]
Choosing $q_A$ to be uniform and using exact relabelling of $G$ makes these
initial scales invariant under bijections.  They remain free up to one
positive multiplier for each cardinality.

For an injection $\iota:B\hookrightarrow A$, define its residual defect by
\[
 K(\iota):=\beta_\iota(q,r)\frac{a_r^B}{a_q^A}.
\]
Equation~\eqref{pt:eq:within-scale} shows that this number is independent of
$q$ and $r$: changing either prior merely inserts and cancels the
corresponding within-set coefficient.  Equation~\eqref{pt:eq:branch-cocycle}
gives
\[
 K(\iota\circ\jmath)=K(\iota)K(\jmath).
\]
Exact relabelling implies that $K(\iota)$ depends only on
$n=|A|$ and $m=|B|$; write it $K_{n,m}$.  Thus
\[
 K_{n,\ell}=K_{n,m}K_{m,\ell},\qquad K_{n,n}=1.
\]
Set $t_n:=K_{n,2}$ and $t_2:=1$.  Then
$K_{n,m}=t_n/t_m$.  Replacing every $a_q^A$ on an $n$-point set by
$t_n a_q^A$ changes the defect to
$K_{n,m}t_m/t_n=1$, while preserving~\eqref{pt:eq:within-scale} and
relabel invariance.  This proves~\eqref{pt:eq:facescale}.

For a boundary prior, $a_r$ means the scale on its support; the preceding
construction says precisely that this agrees with the scale obtained by
reaching that support as a face.  Divide~\eqref{pt:eq:branchaggregation} by
$a_q$ and use~\eqref{pt:eq:facescale}.  Every nondegenerate branch coefficient
becomes
\[
 \frac1{a_q}\,\beta_A(q,r_i)G_{r_i}
 =\frac1{a_{r_i}}G_{r_i}=V_{r_i}.
\]
Singleton values vanish, so their scales and coefficients may be fixed
arbitrarily.  This proves~\eqref{pt:eq:chainrule}.
\end{proof}

\paragraph{Compatibility of the product and branch calibrations.}

The product and sequential evaluations of full revelation now meet.  Their
compatibility first removes the provisional product interaction and then
leaves a strongly additive uncertainty functional.

\begin{lemma}[Interaction collapse and universal chain scale]
\label{pt:lem:scalecoherence}
The coefficient $\kappa$ in~\eqref{pt:eq:QA} is zero.  Moreover there is a
single $C>0$ such that $a_q=C$ for every finite prior with support of size at
least two, and the same value may be assigned on singleton supports.
Consequently the final representatives $V_q=G_q/C$ satisfy, for all nonempty
finite sets $A,B$, all priors $p\in\Delta(A)$ and $r\in\Delta(B)$, and all
$\mu\in\mathcal M_p$ and $\nu\in\mathcal M_r$,
\begin{equation}
 V_{p\otimes r}(\mu\otimes\nu)
   =V_p(\mu)+V_r(\nu).
\tag{PA}\label{pt:eq:PArecovered}
\end{equation}
The chain rule~\eqref{pt:eq:final-chain} extends from full-support fibres to
all finite priors by exact support transport.
The cross-prior bridge of Lemma~\ref{pt:lem:blockbridge} is valid with $V$
in place of $G$.
\end{lemma}

\begin{proof}
In this proof only, put
\[
 H_G(q):=G_q(\chi_q),\qquad Z(q):=1+\kappa H_G(q).
\]
For nondegenerate $q$, Axiom~\textup{(A4)} and zero normalisation give
$H_G(q)>0$, while the strict slice-slope
\eqref{pt:eq:QAslope} gives
\begin{equation}
 Z(q)>0.
\label{pt:eq:Zpositive}
\end{equation}
For a singleton, $H_G=0$ and $Z=1$.

\emph{Product revelation versus sequential revelation.}
For nondegenerate full-support $p,r$, product quasi-additivity gives
\begin{equation}
 H_G(p\otimes r)
 =H_G(p)+H_G(r)+\kappa H_G(p)H_G(r).
\tag{QH}\label{pt:eq:QH}
\end{equation}
Compute the same full revelation by first revealing the $A$-coordinate and
then the $B$-coordinate.  The first-stage law is
$\chi_p\otimes\delta_r$ and has value $H_G(p)$ by~\eqref{pt:eq:QA}.
After action $a$ is revealed, the posterior is $e_a\otimes r$ on the face
$\{a\}\times B$; support restriction and relabelling identify the
continuation value with $H_G(r)$.  Every such branch has coefficient
$a_{p\otimes r}/a_r$ by~\eqref{pt:eq:facescale}.  Since the branch weights
sum to one, Lemma~\ref{pt:lem:branchagg} gives
\begin{equation}
 H_G(p\otimes r)=H_G(p)+\frac{a_{p\otimes r}}{a_r}H_G(r).
\tag{SA}\label{pt:eq:seqA}
\end{equation}
Comparing~\eqref{pt:eq:QH} and~\eqref{pt:eq:seqA}, and using
$H_G(r)>0$, yields
\[
 \frac{a_{p\otimes r}}{a_r}=Z(p).
\tag{R1}\label{pt:eq:ratio1}
\]
Revealing $B$ first gives symmetrically
\[
 \frac{a_{p\otimes r}}{a_p}=Z(r).
\tag{R2}\label{pt:eq:ratio2}
\]
Thus $a_rZ(p)=a_pZ(r)$ for every pair of nondegenerate priors, even on
different finite sets.  By~\eqref{pt:eq:Zpositive}, there is a universal
$C>0$ such that
\begin{equation}
 a_q=CZ(q).
\tag{AS}\label{pt:eq:ascaleZ}
\end{equation}
Assign $a=C$ on singleton supports.

\emph{The grouping weight.}
Let $q$ have full support and let
$\mathcal P=\{A_1,\ldots,A_m\}$ partition its action set into nonempty
cells.  Put
\[
 p_i=q(A_i),\qquad q^i=q(\,\cdot\mid A_i),\qquad
 p=(p_1,\ldots,p_m).
\]
If $C_{\mathcal P}$ reports only the cell, undetectable-distinctions
neutrality gives
\[
 (q,C_{\mathcal P})\sim(p,\Id_{\{1,\ldots,m\}}).
\]
The cross-prior bridge therefore identifies the coarse value as
\begin{equation}
 G_q(\mu_{q,C_{\mathcal P}})=H_G(p).
\tag{BG}\label{pt:eq:blockvalue}
\end{equation}
Append full revelation within each cell.  The final law is $\chi_q$, and the
branch formula together with~\eqref{pt:eq:ascaleZ} gives
\begin{equation}
 H_G(q)=H_G(p)+Z(q)\sum_i p_i\frac{H_G(q^i)}{Z(q^i)}.
\tag{GR}\label{pt:eq:groupingH}
\end{equation}
Here singleton cells contribute zero.

If $\kappa=0$, then $Z\equiv1$ and the conclusion follows.  Otherwise set
$w(q):=1/Z(q)>0$.  Because $\kappa\ne0$ in this case,
$\kappa H_G(x)=Z(x)-1$.  Substitution into
\eqref{pt:eq:groupingH} gives
\[
 Z(q)-1=Z(p)-1+
 Z(q)\sum_i p_i\left(1-\frac1{Z(q^i)}\right).
\]
Cancelling and inverting yields the positive grouping-weight equation
\begin{equation}
 w(q)=w(p)\sum_i p_iw(q^i).
\tag{W}\label{pt:eq:weight}
\end{equation}
All distributions in what follows are read on their positive-probability
supports.  For $u\otimes v$, partition by the first coordinate: the coarse
distribution is $u$ and every conditional distribution is $v$.  Thus
\eqref{pt:eq:weight} gives
\begin{equation}
 w(u\otimes v)=w(u)w(v).
\tag{M}\label{pt:eq:wmult}
\end{equation}

\emph{Two groupings force the weight to be constant.}
Take arbitrary finite probability vectors $u,v$ and write
$x=w(u)$ and $y=w(v)$.  Let $T$ be the law of a triple
$(\ell,z_1,z_2)$: $\ell$ is uniform on $\{0,1\}$, and conditional on
$\ell=0$ the pair $(z_1,z_2)$ has law $u\otimes u$, while conditional on
$\ell=1$ it has law $v\otimes v$.  In tagged notation,
\[
 T:=\tfrac12(u\otimes u)^0\sqcup\tfrac12(v\otimes v)^1.
\]
Grouping first by $\ell$ and using~\eqref{pt:eq:wmult} gives
\begin{equation}
 w(T)=w(\tfrac12,\tfrac12)\frac{x^2+y^2}{2}.
\tag{E1}\label{pt:eq:twoeval1}
\end{equation}
Instead group first by $(\ell,z_1)$.  That marginal is
$S=\tfrac12u^0\sqcup\tfrac12v^1$, so
\[
 w(S)=w(\tfrac12,\tfrac12)\frac{x+y}{2}.
\]
The conditional law of $z_2$ is $u$ when $\ell=0$ and $v$ when $\ell=1$.
A second use of~\eqref{pt:eq:weight} therefore gives
\begin{equation}
 w(T)=w(S)\frac{x+y}{2}
 =w(\tfrac12,\tfrac12)\left(\frac{x+y}{2}\right)^2.
\tag{E2}\label{pt:eq:twoeval2}
\end{equation}
The common leading factor is positive.  Comparing
\eqref{pt:eq:twoeval1} and~\eqref{pt:eq:twoeval2} yields $(x-y)^2=0$.
Thus $w$ is constant on all finite probability vectors.  Its singleton value
is one, so $w\equiv1$.

It follows that $Z\equiv1$, hence $\kappa H_G(q)=0$ for every $q$.
Axiom~\textup{(A4)} supplies a nondegenerate $q$ with $H_G(q)>0$, so
$\kappa=0$, contradicting the temporary alternative.  Equation
\eqref{pt:eq:ascaleZ} now gives $a_q=C$ universally.
Dividing by $C$ proves \eqref{pt:eq:PArecovered} and
\eqref{pt:eq:final-chain} at full support; exact support transport extends
them to boundary priors.  The common positive division also preserves the
cross-prior bridge.
\end{proof}

\subsubsection{Entropy identification}

\begin{lemma}[Entropy reduction and Faddeev identification]
\label{pt:lem:faddeevsketch}
For every nonempty finite $A$ and every $q\in\Delta(A)$, define the value of
full revelation by
\[
 \mathcal H(q):=V_q(\chi_q),
\]
with boundary priors read on their support.
For every nonempty finite $X$ and every channel $P:A\to\Delta(X)$, put
$m_x:=\sum_aq(a)P(x\mid a)$ and, for each reached record $x$,
$r_x(a):=q(a)P(x\mid a)/m_x$.  Then
\begin{equation}
 V_q(\mu_{qP})
 =\mathcal H(q)-\sum_{x:m_x>0}m_x\mathcal H(r_x).
\tag{ER}\label{pt:eq:entropy-reduction}
\end{equation}
Moreover, there is a single $\alpha>0$, independent of all finite alphabets
and priors, such that
\[
 \mathcal H(q)=\alpha\Sh(q)
\]
and consequently
\[
 V_q(\mu_{qP})=\alpha I_{qP}(A;X).
\tag{MI}\label{pt:eq:value-MI}
\]
\end{lemma}

\begin{proof}
\emph{Entropy reduction.}
Assume first that $q$ has full support.  Append full revelation after $P$.
The final posterior law is $\chi_q$, while the continuation after record $x$
has full-revelation law $\chi_{r_x}$ on the support of $r_x$.  The exact
chain rule~\eqref{pt:eq:final-chain} therefore gives
\[
 \mathcal H(q)=V_q(\mu_{qP})
 +\sum_{x:m_x>0}m_x\mathcal H(r_x),
\]
which is~\eqref{pt:eq:entropy-reduction}.  A boundary prior is reduced to its
support; null action rows and null records affect neither side.

\emph{Strong additivity.}
Let $\mathcal P=\{A_1,\ldots,A_m\}$ be a partition of the support of $q$,
and put
\[
 p_i=q(A_i),\qquad q^i=q(\,\cdot\mid A_i),\qquad
 p=(p_1,\ldots,p_m).
\]
In~\eqref{pt:eq:groupingH}, Lemma~\ref{pt:lem:scalecoherence} gives
$Z\equiv1$ and $H_G=C\mathcal H$.  Dividing that identity by $C$ yields
\begin{equation}
 \mathcal H(q)=\mathcal H(p)+\sum_i p_i\mathcal H(q^i).
\tag{FA}\label{pt:eq:faddeev-recursion}
\end{equation}
Relabelling invariance makes $\mathcal H$ symmetric, support restriction makes
it expansive, and $\mathcal H$ vanishes on point masses.  Record processing
gives $\mathcal H\ge0$, while Axiom~\textup{(A4)} gives strict positivity at
nondegenerate full-support priors.  Equation~\eqref{pt:eq:faddeev-recursion}
also allows zero outer weights by deleting and reinserting null cells.

\emph{Continuity: the binary case.}
Write $h(t):=\mathcal H(t,1-t)$.  We derive, rather than assume, its
continuity from Axiom~\textup{(A2)}.  Fix $c\ge0$.  Choose a full-support
binary prior $\theta$ with $\mathcal H(\theta)>0$, as supplied by
Axiom~\textup{(A4)}.  Choose $n\ge1$ and $\lambda\in[0,1]$ so that
\[
 c=\lambda n\mathcal H(\theta).
\]
For $\theta_n:=\theta^{\otimes n}$, product additivity
\eqref{pt:eq:PArecovered} gives
$\mathcal H(\theta_n)=n\mathcal H(\theta)$.  Let $E_{n,\lambda}$ reveal the
$n$-bit action fully with probability $\lambda$ and report a common erasure
symbol otherwise.  Its posterior law is
$\lambda\chi_{\theta_n}+(1-\lambda)\delta_{\theta_n}$, so affinity gives
\[
 V_{\theta_n}(\mu_{\theta_n,E_{n,\lambda}})=c.
\]

Now fix the block environment
\[
 K_c:=\Id_{\{0,1\}}\sqcup E_{n,\lambda}.
\]
Let $q_t^0$ put the binary prior $(t,1-t)$ on its first block and let
$\theta_n^1$ put $\theta_n$ on its second block.  The cross-prior bridge says
\begin{align*}
 h(t)\ge c&\quad\Longleftrightarrow\quad q_t^0\succeq_{K_c}\theta_n^1,\\
 h(t)\le c&\quad\Longleftrightarrow\quad \theta_n^1\succeq_{K_c}q_t^0.
\end{align*}
The environment is fixed and $t\mapsto q_t^0$ is continuous, so
Axiom~\textup{(A2)} makes both level sets closed for every $c\ge0$.  For a
threshold $c<0$, the superlevel set is all of $[0,1]$ and the sublevel set is
empty because $h\ge0$.
Thus every upper and lower level set of $h$ is closed, and $h$ is continuous.

\emph{Continuity on every finite simplex.}
Proceed by induction on the number of coordinates.  For
$q=(q_1,(1-q_1)\bar q)$ with $q_1<1$, recursion
\eqref{pt:eq:faddeev-recursion} for the partition
$\{\{1\},\{2,\ldots,n\}\}$ gives
\begin{equation}
 \mathcal H(q)=h(q_1)+(1-q_1)\mathcal H(\bar q).
\label{pt:eq:continuity-induction}
\end{equation}
The right side is continuous away from $q_1=1$ by binary continuity and the
induction hypothesis.  At $q_1=0$, support restriction and $h(0)=0$ give
the same formula on the face.  At $q_1=1$, the induction hypothesis makes
$\mathcal H$ bounded on the compact $(n-1)$-simplex, so the second term in
\eqref{pt:eq:continuity-induction} tends to zero and the first tends to
$h(1)=0$.  Relabelling handles the other faces.  Hence $\mathcal H$ is
continuous on each closed finite simplex.

We have now verified all hypotheses of the strong-additivity form of
Faddeev's theorem.  Axiom~\textup{(A2)} supplied continuity;
Lemma~\ref{pt:lem:transport} supplied relabelling invariance and
expansibility; $\mathcal H$ vanishes on point masses; and
Lemmas~\ref{pt:lem:chain} and~\ref{pt:lem:scalecoherence} supplied strong
additivity~\eqref{pt:eq:faddeev-recursion}.  Record processing has already
established $\mathcal H\ge0$.  The finite entropy characterisation
\cite[Theorem~6]{BaezFritzLeinster2011} therefore gives
\[
 \mathcal H(q)=\alpha\Sh(q)
\]
for some $\alpha\ge0$.  Axiom~\textup{(A4)} forces $\alpha>0$.
Substitution in~\eqref{pt:eq:entropy-reduction} gives
\[
 V_q(\mu_{qP})
 =\alpha\left(\Sh(q)-\sum_x m_x\Sh(r_x)\right)
 =\alpha I_{qP}(A;X),
\]
where the sum is over positive-mass records.  This proves
\eqref{pt:eq:value-MI}.
\end{proof}

\paragraph{Conclusion.}
The cross-prior bridge and Lemma~\ref{pt:lem:faddeevsketch} give, for arbitrary
pairs $(q,P)$ and $(r,Q)$,
\[
 q^0\succeq_{P\sqcup Q}r^1
 \quad\Longleftrightarrow\quad
 I_{qP}\ge I_{rQ},
\tag{BMI}\label{pt:eq:final-block-MI}
\]
including boundary priors by support restriction.  With $Q=P$, the
duplication clause of Axiom~\textup{(A5)} gives the proposition's fixed-channel
claim.  The positive multiplier $\alpha$ cancels from these comparisons; its
cardinal calibration against payoff utility is the next step.  Thus the
conclusion here is ordinal: it fixes the common ranking of pure-trace pairs,
not their exchange rate against material utility.  For distinct
$i,j\in J$, the irrelevant-block clause of Axiom~\textup{(A5)} reduces an
arbitrary block union to the canonical two-block comparison:
\[
 q_i^i\succeq_{\bigsqcup_{k\in J}P_k}q_j^j
 \quad\Longleftrightarrow\quad
 q_i^0\succeq_{P_i\sqcup P_j}q_j^1.
\]
On the information side, the block tag is constant under a block-supported
prior, so
\[
 I_{q_i^i,\,\bigsqcup_{k\in J}P_k}=I_{q_iP_i},
\]
and similarly for block $j$.  Applying~\eqref{pt:eq:final-block-MI} to the
right-hand side proves the block claim and completes
Proposition~\ref{pt:prop:main}.
